\documentclass[../DoAn.tex]{subfiles}
\begin{document}

\begin{center}
    \Large{\textbf{TÓM TẮT NỘI DUNG ĐỒ ÁN}}\\
\end{center}
\vspace{1cm}
Đồ án tập trung vào việc nghiên cứu và xây dựng hệ thống RAG (Retrieval-Augmented Generation) Chatbot hỗ trợ sinh viên hỏi-đáp tự động về các quy định, quy chế đào tạo của Đại học Bách khoa Hà Nội. Mục tiêu chính là cung cấp một công cụ tra cứu thông minh, cho phép sinh viên đặt câu hỏi bằng ngôn ngữ tự nhiên và nhận được câu trả lời chính xác kèm theo trích dẫn nguồn cụ thể từ các văn bản quy chế.

Hệ thống được xây dựng theo kiến trúc RAG hoàn chỉnh, bao gồm các thành phần chính: xử lý PDF đa phương thức (Docling, PyMuPDF4LLM, olmOCR), chunking văn bản pháp lý theo cấu trúc Điều-Khoản, embedding với mô hình E5 Multilingual (1024 chiều), lưu trữ vector sử dụng FAISS, và sinh câu trả lời bằng Gemini 2.5 Flash. Backend được phát triển với FastAPI và frontend sử dụng React 18 với TypeScript.

Dữ liệu đầu vào bao gồm 16 văn bản quy định đào tạo của ĐHBK Hà Nội, sau khi xử lý tạo ra 1,247 chunks được lưu trong vector store. Hệ thống đạt độ chính xác xử lý tiếng Việt 99.5\%+ và thời gian phản hồi 2-4 giây cho mỗi câu hỏi.

Đồ án là nền tảng quan trọng cho việc ứng dụng công nghệ RAG trong môi trường giáo dục, hỗ trợ sinh viên tra cứu thông tin quy chế một cách hiệu quả và chính xác.

\begin{flushright}
Sinh viên thực hiện\\
\begin{tabular}{@{}c@{}}
\textit{(Ký và ghi rõ họ tên)}
\end{tabular}
\end{flushright}

\end{document}
